% =============================================================================
% Machine Learning & Deep Learning -- Project Report Template
% Alpen-Adria-Universität Klagenfurt
%
% Page limit: 9 pages, EXCLUDING references. Everything else counts,
% including figures, tables, the appendix and the AI disclosure statement.
% Do not change margins, font sizes or spacing.
%
% The comments under each heading describe what the section must contain.
% Replace every "Your text." with your content and submit the PDF in Moodle,
% together with your code (see Section 5 of the project guide, README.md).
% =============================================================================

\documentclass{article}

% "preprint" shows your names; without it the template produces an
% anonymous submission.
\usepackage[preprint]{neurips_2025}

\usepackage[utf8]{inputenc} % allow utf-8 input
\usepackage[T1]{fontenc}    % use 8-bit T1 fonts
\usepackage{hyperref}       % hyperlinks
\usepackage{url}            % simple URL typesetting
\usepackage{booktabs}       % professional-quality tables
\usepackage{amsmath}        % mathematical environments
\usepackage{amsfonts}       % blackboard math symbols
\usepackage{nicefrac}       % compact symbols for 1/2, etc.
\usepackage{microtype}      % microtypography
\usepackage{xcolor}         % colors
\usepackage{graphicx}       % figures

\title{Your Project Title}

\author{%
  First Student \\
  Student ID: 00000000 \\
  \texttt{first.student@edu.aau.at} \\
  \And
  Second Student \\
  Student ID: 00000000 \\
  \texttt{second.student@edu.aau.at} \\
  \And
  Third Student \\
  Student ID: 00000000 \\
  \texttt{third.student@edu.aau.at} \\
}

\begin{document}

\maketitle

\begin{abstract}
% One paragraph summarizing the problem, your approach and your main result.
Your abstract.
\end{abstract}

% =============================================================================
\section{Introduction}
% - The problem you are solving and why it is important.
% - Your hypothesis and what you contribute.
% - A short summary of your results.

Your text.

% =============================================================================
\section{Related Work}
% - Which models and approaches from the literature suit this problem and your
%   data, and why?
% - How does your work relate to them?
% - Cite all sources properly, e.g. \citet{bishop2024deep} or
%   \citep{goodfellow2016deep}, and add them to references.bib.

Your text.

% =============================================================================
\section{Data}
% The data you use, with examples:
% - Where does it come from?
% - What are its properties, such as class balance, missing values and scale,
%   and how do you deal with them?
% - Which preprocessing, filtering, augmentation or other manipulations do you
%   apply?
% - How do you split the data into training, validation and test sets, and how
%   do you prevent information from the test set leaking into training?

Your text.

% =============================================================================
\section{Method}
% Your approach and why it fits the problem. In a group, describe the approach
% of every member and state who did what.
% - Model architecture, output activation and loss function.
% - Training procedure: optimizer, learning-rate schedule, regularization.
% - Which hyperparameters does your approach have, and how can they influence
%   the results?
% - How do these choices relate to your hypothesis?
% - You must show that you applied methods covered during the semester to your
%   problem. Figures of your architecture or pipeline are welcome.

Your text.

% =============================================================================
\section{Experimental Setup}
% - Baselines: at least one simple baseline, and one from the literature,
%   blogs, Kaggle, etc., if available.
% - Fair comparison: the same data split, test data, evaluation measures,
%   random seeds and a comparable training budget for all approaches.
% - Ablation design: which components each member removes or replaces.
% - Evaluation measures and why they fit the task.
% - Hyperparameter search: search space, strategy, and the data used for model
%   selection.
% - Implementation details: framework, hardware and training time.
% - Links to your code (archive, repository or Colab notebook) and to your
%   public W&B project (see the paragraph below).

Your text.

\paragraph{Code and experiments.}
Code: \url{https://github.com/your-group/your-project} (or a Colab notebook). \\
Experiment tracking: \url{https://wandb.ai/your-entity/your-project}.

% =============================================================================
\section{Results}

\subsection{Comparative Analysis}
% REQUIRED: the comparative analysis of all approaches is the main result of the
% project. A report without it is incomplete.
% - Compare the approaches of all group members with each other and with the
%   baselines from the Experimental Setup: same data split, test data, evaluation
%   measures and random seeds, and a comparable training budget.
% - Present the main results in one table (see Table~\ref{tab:results}): the
%   main measures (mean and standard deviation over several random seeds), the
%   number of parameters and the training time. Add learning curves where they
%   help (see Figure~\ref{fig:curves}).
% - Does each approach beat the baselines, and by how much?
% - Explain the differences: which approach works best, why, and at what cost?
%   End with a justified recommendation, e.g. "Is the ResNet's improvement worth
%   its size and training time?"
% - Results and figures should come from runs logged in your W&B project.

Your text.

\begin{table}[t]
  \caption{Test results of all approaches and baselines (mean $\pm$ standard deviation over 3 random seeds), with model size and training time.}
  \label{tab:results}
  \centering
  \small
  \setlength{\tabcolsep}{4pt}
  \begin{tabular}{llccccc}
    \toprule
    Approach & Member & Precision & Recall & ROC-AUC & Parameters & Training time \\
    \midrule
    Simple baseline          & ---    & 0.00 & 0.00 & 0.00 & --- & --- \\
    Baseline from literature & ---    & 0.00 & 0.00 & 0.00 & --- & --- \\
    \midrule
    Approach A               & First  & 0.00 $\pm$ 0.00 & 0.00 $\pm$ 0.00 & 0.00 $\pm$ 0.00 & 0.0\,M & 0\,min \\
    Approach B               & Second & 0.00 $\pm$ 0.00 & 0.00 $\pm$ 0.00 & 0.00 $\pm$ 0.00 & 0.0\,M & 0\,min \\
    Approach C               & Third  & 0.00 $\pm$ 0.00 & 0.00 $\pm$ 0.00 & 0.00 $\pm$ 0.00 & 0.0\,M & 0\,min \\
    \bottomrule
  \end{tabular}
\end{table}

\begin{figure}[t]
  \centering
  % Replace the box with your figure, e.g.:
  % \includegraphics[width=0.8\linewidth]{figures/learning_curves.pdf}
  \fbox{\parbox[c][4cm][c]{0.8\linewidth}{\centering Learning curves (training and validation loss per epoch)}}
  \caption{Training and validation loss of all approaches.}
  \label{fig:curves}
\end{figure}

% =============================================================================
\section{Analysis}

\subsection{Ablation Studies}
% REQUIRED: every group member reports the ablation study of their own approach.
% - Which components of your approach actually contribute to the result?
% - Remove or replace one component at a time, keep everything else fixed, and
%   compare with the full approach (see Table~\ref{tab:ablation}).
% - Did the results confirm the hypotheses from your proposal? Connect each
%   finding to the lecture, e.g. "without skip connections, the deeper network
%   trains worse because of vanishing gradients".
% - State who ran which study.

Your text.

\begin{table}[t]
  \caption{Ablation studies of all approaches: test ROC-AUC (mean $\pm$ standard deviation over 3 random seeds) and change compared with the full approach.}
  \label{tab:ablation}
  \centering
  \small
  \begin{tabular}{llcc}
    \toprule
    Approach (member) & Variant & ROC-AUC & Change \\
    \midrule
    A: plain CNN (First)        & full approach                    & 0.00 $\pm$ 0.00 & --- \\
                                & without batch normalization      & 0.00 $\pm$ 0.00 & $-0.00$ \\
                                & without data augmentation        & 0.00 $\pm$ 0.00 & $-0.00$ \\
    \midrule
    B: ResNet-18 (Second)       & full approach                    & 0.00 $\pm$ 0.00 & --- \\
                                & without skip connections         & 0.00 $\pm$ 0.00 & $-0.00$ \\
                                & random init.\ instead of pre-training & 0.00 $\pm$ 0.00 & $-0.00$ \\
    \midrule
    C: training procedure (Third) & best recipe                    & 0.00 $\pm$ 0.00 & --- \\
                                & SGD instead of Adam              & 0.00 $\pm$ 0.00 & $-0.00$ \\
                                & constant learning rate           & 0.00 $\pm$ 0.00 & $-0.00$ \\
    \bottomrule
  \end{tabular}
\end{table}

\subsection{Error Analysis}
% - On which examples does each approach fail, and why?
% - Which examples do all approaches get wrong? Show typical failure cases.

Your text.

% =============================================================================
\section{Limitations and Conclusions}
% - What are the limits of your results?
% - How do you assess your work, and what would you recommend as next steps?

Your text.

% =============================================================================
\section*{Author Contributions}
% State what each group member has done, e.g. "(Name) implemented and evaluated
% approach A, ...".

Your text.

% =============================================================================
\bibliographystyle{plainnat}
\bibliography{references}

% =============================================================================
\appendix

\section{Additional Results}
% Optional: additional results, the full hyperparameter configuration, and
% interesting code snippets or algorithms. The appendix counts towards the
% page limit.

Your text.

\section{AI Disclosure Statement}
% - Declare any use of generative AI tools and coding assistants.
% - State which tools you used and for what purpose, e.g. brainstorming,
%   literature search, code generation or text drafting.
% - If you did not use any, say so.

Your text.

\end{document}
